\scanpage{31}
\mspage{15}
\unsure{remontées|remonter} en des hypothèses dérivées \inserted{de nature} arithmético-géométrique, on peut laisser tomber les aspects anabéliens, pour jouer à fond sur les aspects abéliens -- mais dans un contexte arithmétique. Au lieu d'avoir \inserted{et d'utiliser} seulement, comme dans le cor. 1 a), des hypothèses que $H_1(\overline f)=H_1(\overline g):H_1(\overline X)\longrightarrow H_1(\overline Y)$, on a une hypothèse plus fine, sur l'égalité de deux homomorphismes \uline{d'extensions} de $E_{\overline K/K}$ par $H_1(\overline X,\widehat{\mathbb Z}_\ell)$ resp. $H_1(\overline Y,\widehat{\mathbb Z}_\ell)$. $\forall\ell\in\mathbb P$.

Démonstration du corollaire 3.
Il suffit de voir que $f(x)=g(x)$ pour tt $x\in X(K)$ (car, quitte à \deleted{étendre} \inserted{faire une extension finie} du corps de base, cela reste vrai \deleted{après tt extension de corps} \inserted{pour $x\in X(\overline K)$}). \deleted{\unsure{faire}}

Or un point définit une classe de \deleted{d'autre part} \inserted{conjugaison de} scindages de $E_{X/K}$ sur $E_{\overline K/K}$, et a fortiori de l'extension à noyau abélien associée -- l'hypothèse implique que $f(x)$ et $g(x)$ définissent la $\unsure{=|\overline m}$ classe de
\margin{30}
